\section{Our Construction}
\label{sec:construction}

In this section we present the formal construction of our verifiable matrix multiplication protocol $\protocol$.
We begin by defining the relation chain that reduces the original matrix multiplication statement to a succinct SNARK-friendly relation (Section~\ref{sec:relation-chain}), then give the full protocol description (Section~\ref{sec:protocol-desc}), and conclude with a complexity analysis (Section~\ref{sec:complexity}).
Table~\ref{tab:notation} summarizes the notation used throughout the paper.

\begin{table}[t]
\centering
\small
\begin{tabular}{@{}cl@{}}
\toprule
\textbf{Symbol} & \textbf{Description} \\
\midrule
\multicolumn{2}{@{}l}{\emph{Parameters}} \\
$\mathbb{F}$                        & Finite field \\
$k$                                  & Matrix dimension ($k \times k$) \\
$n$                                  & Codeword length of $\mathcal{C}$ \\
$\delta$                             & Relative minimum distance of $\mathcal{C}$ \\
$t = |I|$                           & Query complexity (number of sampled indices) \\
$\lambda$                            & Security parameter \\
\midrule
\multicolumn{2}{@{}l}{\emph{Matrices and Vectors}} \\
$\mat{A}, \mat{B}, \mat{C}$         & Input/output matrices in $\mathbb{F}^{k \times k}$ \\
$\mat{D}$                           & Error matrix $\mat{A}\mat{B} - \mat{C}$ \\
$\vect{r}$                          & Challenge vector $(1,\, r,\, r^2, \ldots, r^{k-1})$ \\
$\vect{x}, \vect{y}, \vect{z}$      & Intermediate vectors: $\vect{r}\mat{A}$,\; $\vect{x}\mat{B}$,\; $\vect{r}\mat{C}$ \\
\midrule
\multicolumn{2}{@{}l}{\emph{Encoding and Commitment}} \\
$\mathcal{C}$                        & Linear error-correcting code $\mathbb{F}^k \to \mathbb{F}^n$ \\
$\enc(\cdot)$                        & Encoding function (row-wise for matrices) \\
$\enc(\mat{M})^{(j)}$               & $j$-th column of $\enc(\mat{M}) \in \mathbb{F}^{k \times n}$ \\
$\fold(\vect{c}, \vect{r})$         & Inner product $\langle \vect{r}, \vect{c} \rangle$ \\
$\merkle(\cdot)$                     & Merkle tree commitment \\
$\merkleO(\cm, j)$                   & Merkle tree opening at index $j$ \\
$\merkleV(\pi, v, \cm)$             & Merkle tree verification \\
$\openV(\pi, v, \cm, j)$            & Position-opening verifier of selected backend \\
$\cm_\chi$                          & Merkle root for committed object $\chi$ \\
$\pi_{\chi,j}$                      & Merkle authentication path for $\chi$ at position $j$ \\
$\sigma_v$                          & Code-commitment backend proof/witness for $\cm_v$ \\
\midrule
\multicolumn{2}{@{}l}{\emph{Sets and tuples}} \\
$\mathcal{M}$                        & Matrix set $\{\mat{A}, \mat{B}, \mat{C}\}$ \\
$\mathcal{W}$                        & Vector set $\{\vect{x}, \vect{y}, \vect{z}\}$ \\
$\mathcal{S}$                        & Full committed-object set $\mathcal{M} \cup \mathcal{W}$ \\
$I$                                  & Ordered tuple of queried column indices ($I \in [n]^t$) \\
\midrule
\multicolumn{2}{@{}l}{\emph{SNARK}} \\
$\stmt$                              & SNARK public statement (instance) \\
$\wit$                               & SNARK witness \\
$\pi$                                & SNARK proof \\
$\snarkP, \snarkV$                   & SNARK prover and verifier \\
$\para$                              & Public parameters \\
\bottomrule
\end{tabular}
\caption{Summary of notation.}
\label{tab:notation}
\end{table}

\subsection{Relation Reduction Chain}
\label{sec:relation-chain}

Our construction proceeds through two successive reductions.
Each reduction preserves completeness exactly and introduces a bounded soundness error that we analyze in Section~\ref{sec:security}.

\paragraph{Input-commitment model.}
Following the committed-oracle viewpoint common in code-based proof systems
such as Brakedown~\cite{golovnev2023brakedown} and Orion~\cite{xie2022orion},
we take the input matrix commitments as certified public inputs.
Concretely, $\cm_A,\cm_B,\cm_C$ are assumed to bind to row-wise encodings
$\enc(\mat{A}),\enc(\mat{B}),\enc(\mat{C})$ of some matrices
$\mat{A},\mat{B},\mat{C} \in \mathbb{F}^{k \times k}$.
Certification means that every accepting opening of $\cm_M$, for
$M \in \{\mat{A},\mat{B},\mat{C}\}$, is a coordinate or column of
$\enc(M)$ for a fixed matrix $M$ determined before the verifier samples the
Freivalds challenge.
This certification may be supplied by the surrounding committed-codeword
system, by a public generation process, or by a separate input-opening proof.
It is a precondition of the multiplication-checking protocol rather than an
online cost charged to $\protocol$.
The goal of $\protocol$ is to verify the multiplication relation among these
committed matrices without arithmetizing the full matrix product.
Throughout this section, $\merkle(\enc(\mat{M}))$ for an encoded matrix means
the Merkle root over the length-$n$ sequence of encoded columns
$(\enc(\mat{M})^{(1)},\ldots,\enc(\mat{M})^{(n)})$.
For an encoded vector, $\merkle(\enc(v))$ commits to the length-$n$ sequence of
coordinates.

\begin{definition}[Code-commitment backend]\label{def:codecom-backend}
Fix a systematic linear code $\mathcal{C}: \mathbb{F}^{k}\to\mathbb{F}^{n}$
and a position-opening layer whose verification predicate is written
$\openV(\pi_j,a,\cm,j)$.
For the Merkle instantiation, $\openV$ is exactly $\merkleV$ with the queried
index included in the leaf domain; for KZG or CPLink-style instantiations it is
the corresponding opening or linking verifier.
A code-commitment backend $\codeCom$ provides a prover relation $\codeComP$ and
a verifier relation $\codeComV$ for vector commitments.
It must satisfy the following two properties.

\emph{Completeness.}
For every $v \in \mathbb{F}^{k}$, if
$\cm_v$ commits to $\enc(v)=\mathcal{C}(v)$ under the chosen opening layer, then
an honest prover can produce proof material
$\sigma_v \gets \codeComP(\cm_v,v,\enc(v))$ accepted by
$\codeComV(\cm_v,v,\sigma_v)$.

\emph{Global codeword binding.}
The backend is codeword-binding with error
$\epsilon_{\mathsf{cc}}(\lambda)$ if, for every $\ppt$ adversary
$\mathcal{A}$,
\[
\Pr[\mathsf{Bad}^{\mathcal{A}}_{\mathsf{cc}}(\lambda)=1]
\le \epsilon_{\mathsf{cc}}(\lambda).
\]
The event $\mathsf{Bad}^{\mathcal{A}}_{\mathsf{cc}}$ occurs if
$\mathcal{A}$ outputs $(\cm_v,v,\sigma_v,j,a,\pi_j)$ such that
$\codeComV(\cm_v,v,\sigma_v)=1$, $\openV(\pi_j,a,\cm_v,j)=1$, and
$a \neq \enc(v)_j$.
Thus, except with probability $\epsilon_{\mathsf{cc}}$, any vector commitment
accepted by $\codeComV$ is globally linked to the full codeword $\enc(v)$:
every accepting opening at every coordinate is forced to equal the
corresponding coordinate of $\enc(v)$, not merely the sampled coordinates.
The separate position-binding error of the opening layer is accounted for by
$\epsilon_{\mathsf{bind}}$ in Theorem~\ref{thm:security}.
The cost of verifying $\codeComV(\cm_v,v,\sigma_v)$ is denoted
$E_{\mathsf{cc}}(k)$.
\end{definition}

Rather than introducing a new commitment primitive, $\protocol$ is stated
relative to this committed-codeword interface, which matches the role played by
encoded committed oracles in Brakedown~\cite{golovnev2023brakedown} and
Orion~\cite{xie2022orion}.
Operationally, the backend can be instantiated by a Brakedown/Orion-style
code-commitment layer, an in-circuit encoder check tied to the committed root,
or a commitment-to-codeword linking proof supplied by the chosen backend.

\paragraph{Backend contract used by the theorem.}
Theorem~\ref{thm:security} uses only the following backend facts:
(i)~the input matrix roots are certified commitments to fixed row-wise
codewords before $r$ is sampled;
(ii)~the intermediate vector roots are sent before $I$ is sampled;
(iii)~an accepting $\codeComV(\cm_v,v,\sigma_v)$ globally links $\cm_v$ to
$\enc(v)$ as in Definition~\ref{def:codecom-backend}; and
(iv)~the opening layer is position-binding for every queried coordinate.
No part of the proof relies on a new algebraic assumption specific to
$\protocol$ beyond this committed-codeword interface.

\paragraph{Backend instantiation boundary.}
The security theorem is compositional: it proves the multiplication-checking
protocol relative to any backend satisfying Definition~\ref{def:codecom-backend}.
An end-to-end system must choose one concrete backend and specify its algorithms,
the contents of $\sigma_v$, and the constraint/proof-size contribution of
$\codeComV$.
Our implementation evaluates the local fold/opening/linking layer under a
linear-time $E_{\mathsf{cc}}(k)=\mathcal{O}(k)$ backend model; it does not yet
implement or measure $\codeComV$ independently.

\paragraph{The original relation $\mathcal{R}_{\mathsf{orig}}$.}
The starting point is the statement that the prover knows matrices whose
\emph{encoded} commitments satisfy the multiplication relation.
The public instance is $(\cm_A, \cm_B, \cm_C)$ and the witness is
$(\mat{A}, \mat{B}, \mat{C})$; the relation $\mathcal{R}_{\mathsf{orig}}$
requires:
\begin{align*}
  &\cm_A = \merkle(\enc(\mat{A})) \;\land\;
    \cm_B = \merkle(\enc(\mat{B})) \;\land\;
    \cm_C = \merkle(\enc(\mat{C})), \\
  &\mat{A}\mat{B} = \mat{C}.
\end{align*}
Thus a matrix commitment binds the verifier to the row-wise codeword matrix
rather than to the raw matrix alone.
Since $\mathcal{C}$ is systematic, the raw matrix is also determined by the
first $k$ coordinates of each committed row.
Proving membership in $\mathcal{R}_{\mathsf{orig}}$ directly inside a SNARK would require $\mathcal{O}(k^3)$ multiplication constraints, which is exactly what we aim to avoid.

\paragraph{Reduction 1: Freivalds' randomization ($\mathcal{R}_{\mathsf{orig}} \rightsquigarrow \mathcal{R}_{\mathsf{Freivalds}}$).}
Using the verifier's random challenge $r \in \mathbb{F}$, we construct the vector $\vect{r} = (1, r, r^2, \ldots, r^{k-1})$ and define intermediate vectors $\vect{x} = \vect{r}\mat{A}$, $\vect{y} = \vect{x}\mat{B}$, $\vect{z} = \vect{r}\mat{C}$.
Let $\mathcal{S} = \{\mat{A}, \mat{B}, \mat{C}, \vect{x}, \vect{y}, \vect{z}\}$ denote the full set of committed objects.
The instance is $(\{\cm_{\chi}\}_{\chi \in \mathcal{S}}, r)$ and the witness is $\{\chi\}_{\chi \in \mathcal{S}}$;
the relation $\mathcal{R}_{\mathsf{Freivalds}}$ requires:
\begin{align*}
  &\textstyle\bigwedge_{\chi \in \mathcal{S}}\; \cm_{\chi} = \merkle(\enc(\chi)), \\
  &\vect{r}\mat{A} = \vect{x} \;\land\; \vect{x}\mat{B} = \vect{y}, \\
  &\vect{r}\mat{C} = \vect{z} \;\land\; \vect{y} = \vect{z}.
\end{align*}
By the Schwartz--Zippel lemma, if $\mat{A}\mat{B} \neq \mat{C}$, then $\vect{y} \neq \vect{z}$ with probability at least $1 - (k-1)/|\mathbb{F}|$.
Note that while this relation eliminates the $\mathcal{O}(k^3)$ matrix product, verifying the vector-matrix products still costs $\mathcal{O}(k^2)$ each.

\paragraph{Reduction 2: Proximity lifting ($\mathcal{R}_{\mathsf{Freivalds}} \rightsquigarrow \mathcal{R}_{\mathsf{SNARK}}$).}
We now encode all matrices and vectors using a systematic $\mathbb{F}$-linear code $\mathcal{C}: \mathbb{F}^k \to \mathbb{F}^n$ with relative distance $\delta$, and replace the global vector-matrix checks with \emph{local} consistency checks at random positions.
Let $I=(j_1,\ldots,j_t) \sample [n]^t$ be the ordered tuple of queried column
indices sampled independently with replacement, and let
$\mathcal{M} = \{\mat{A}, \mat{B}, \mat{C}\}$,
$\mathcal{W} = \{\vect{x}, \vect{y}, \vect{z}\}$ denote the matrix set and
vector set respectively.

The \emph{public instance} is
$\stmt = \bigl(\{\cm_{\chi}\}_{\chi \in \mathcal{S}},\; r,\; I\bigr)$
and the \emph{witness} is
\[
  \wit = \Bigl(\vect{x},\, \vect{y},\, \vect{z},\;
    \sigma_x,\,\sigma_y,\,\sigma_z,\;
    \bigl\{\enc_{\chi}^{(j)},\, \pi_{\chi,j}\bigr\}_{j \in I,\, \chi \in \mathcal{S}}\Bigr).
\]
The relation $\mathcal{R}_{\mathsf{SNARK}}$ requires $(\stmt, \wit)$ to satisfy all of:
\begin{align}
  &\codeComV(\cm_x,\vect{x},\sigma_x)=1
    \;\land\; \codeComV(\cm_y,\vect{y},\sigma_y)=1 \notag\\
  &\quad\;\land\; \codeComV(\cm_z,\vect{z},\sigma_z)=1,
    \tag{CODECOM}\label{eq:codecom-check}\\
  &\forall j \in I{:}\quad
    \enc_{x,j} = \fold\!\bigl(\enc_{A}^{(j)},\,\vect{r}\bigr), \tag{FOLD-A}\label{eq:fold-A}\\
  &\forall j \in I{:}\quad
    \enc_{y,j} = \fold\!\bigl(\enc_{B}^{(j)},\,\vect{x}\bigr), \tag{FOLD-B}\label{eq:fold-B}\\
  &\forall j \in I{:}\quad
    \enc_{z,j} = \fold\!\bigl(\enc_{C}^{(j)},\,\vect{r}\bigr), \tag{FOLD-C}\label{eq:fold-C}\\
  &\forall j \in I{:}\quad
    \enc_{y,j} = \enc_{z,j}, \tag{EQ}\label{eq:eq-check}\\
  &\forall \chi \in \mathcal{S},\, j \in I{:}\quad
    \merkleV\!\bigl(\pi_{\chi,j},\,\enc_{\chi}^{(j)},\,\cm_{\chi}\bigr) = 1.
    \tag{PATH}\label{eq:path-check}
\end{align}
Here $\forall j \in I$ ranges over every sampled position in the ordered tuple,
including repetitions.
For a matrix commitment, $\enc_{\chi}^{(j)}$ denotes the full $j$-th encoded
column; for a vector commitment, it denotes the scalar coordinate
$\enc(\chi)_j$.
Equation~\eqref{eq:path-check} is written for the Merkle instantiation used in
the transparent variant.
In KZG or CPLink-style variants, the same slot is occupied by the selected
opening or linking verifier; the security proof uses only position binding of
the opening layer and the codeword-binding property of
Definition~\ref{def:codecom-backend}.

The critical property is that $\mathcal{R}_{\mathsf{SNARK}}$ involves only $\mathcal{O}(|I| \cdot k)$ constraints for the fold operations, plus $\mathcal{O}(|I| \cdot \log n)$ hash checks for the Merkle paths and the backend code-commitment cost $E_{\mathsf{cc}}(k)$.
Since $|I|$ is a security-parameter-sized constant independent of $k$, the circuit size is \emph{linear in $k$}.

\paragraph{Codeword-validity cost.}
The backend checks in~\eqref{eq:codecom-check} force the prover to commit to
\emph{full} valid vector codewords before the query tuple is sampled.
Without this binding step, a cheating prover could choose sampled openings that
satisfy the fold checks without committing to any genuine encoding of a fixed
vector, defeating the distance-based soundness argument.
Analogously, the input matrix commitments $\cm_A,\cm_B,\cm_C$ must be
well-formed commitments to row-wise encodings of $\mat{A},\mat{B},\mat{C}$.
In our model this is part of the public input-commitment interface, not a
property re-proved by the multiplication protocol itself.
For a general linear code represented by a dense generator matrix
$G \in \mathbb{F}^{k \times n}$, the constraint $\enc_v = v \cdot G$
requires $\mathcal{O}(k(n-k))$ inner-product work.
For linear-time encodable codes (e.g., expander-based codes~\cite{golovnev2023brakedown})
and a code-commitment backend whose verifier follows the encoder circuit, the
same validity check can be arithmetized with $\mathcal{O}(n)$
constraints, which is $\mathcal{O}(k)$ at constant rate.
Thus the backend codeword-validity term is linear for the code family used in
Table~\ref{tab:meow-vs-freivalds}, but becomes $\mathcal{O}(k^2)$ for dense
general constant-rate codes.
We discuss the trade-off in Section~\ref{sec:complexity}.

\subsection{Protocol Description}
\label{sec:protocol-desc}

We now give the full specification of $\protocol$.
Let $\mathcal{C}: \mathbb{F}^k \to \mathbb{F}^n$ be a systematic linear code with relative distance $\delta$ and $\merkle$ denote the Merkle tree commitment.
Let $\mathcal{M} = \{\mat{A}, \mat{B}, \mat{C}\}$ be the matrix set, $\mathcal{W} = \{\vect{x}, \vect{y}, \vect{z}\}$ the vector set, and $\mathcal{S} = \mathcal{M} \cup \mathcal{W}$.
The verifier is given certified input commitments
$(\cm_A,\cm_B,\cm_C)$ to $\enc(\mat{A}),\enc(\mat{B}),\enc(\mat{C})$.
The protocol $\protocol$ operates between a prover $\mathcal{P}$ and verifier $\mathcal{V}$ as shown in Figure~\ref{fig:protocol}.

\begin{figure*}[ht!]
\centering
\scalebox{0.85}{%
    \fbox{%
    \procedure[]{\protocol}{%
    \mathcal{P} \< \< \mathcal{V} \\[4pt]
    %% Phase 1
    \textbf{Phase 1: Public input} \< \< \\
    \text{Certified roots } \cm_A,\cm_B,\cm_C
      \text{ bind to } \enc(\mat{A}),\enc(\mat{B}),\enc(\mat{C}) \< \< \\[4pt]
    %% Phase 2
    \< \< \textbf{Phase 2: Challenge} \\
    \< \sendmessageleft*[2.5cm]{r} \< r \sample \FF \\[4pt]
    %% Phase 3
    \textbf{Phase 3: Respond} \< \< \\
    \vect{r} \gets (1, r, r^2, \dots, r^{k-1}) \< \< \\
    \vect{x} \gets \vect{r}\mat{A},\quad \vect{y} \gets \vect{x}\mat{B},\quad \vect{z} \gets \vect{r}\mat{C} \< \< \\
    \text{For } v \in \mathcal{W}: \< \< \\
    \quad \enc_v \gets \enc(v),\quad \cm_v \gets \merkle(\enc_v) \< \< \\
    \quad \sigma_v \gets \codeComP(\cm_v,v,\enc_v)
      \< \sendmessageright*[2.5cm]{\cm_x,\cm_y,\cm_z} \< \\[4pt]
    %% Phase 4
    \< \< \textbf{Phase 4: Query} \\
    \< \sendmessageleft*[2.5cm]{I} \< I=(j_1,\ldots,j_t) \sample [n]^{t} \\[4pt]
    %% Phase 5
    \textbf{Phase 5: Open \& Prove} \< \< \\
    \text{For } j \in I,\ \chi \in \mathcal{S}: \< \< \\
    \quad (\enc_\chi^{(j)}, \pi_{\chi,j}) \gets \merkleO(\cm_\chi, j) \< \< \\
    \stmt \gets \Big( \{\cm_{\chi}\}_{\chi \in \mathcal{S}},\, r,\, I \Big) \< \< \\
    \wit \gets \Big( \vect{x}, \vect{y}, \vect{z},\, \sigma_x,\sigma_y,\sigma_z,\,
      \{\enc_{\chi}^{(j)}, \pi_{\chi,j}\}_{j \in I, \chi \in \mathcal{S}} \Big) \< \< \\
    \pi \gets \snarkP(\para, \stmt, \wit) \< \sendmessageright*[2.5cm]{\pi} \< \\[4pt]
    %% Phase 6
    \< \< \textbf{Phase 6: Verify} \\
    \< \< b \gets \snarkV(\para, \stmt, \pi) \\
    \< \< \text{Accept iff } b = 1
    }%
    }%
}
\caption{The $\protocol$ protocol for verifiable matrix multiplication.
$\mathcal{S} = \mathcal{M} \cup \mathcal{W}$ is the full set of committed objects, $t = |I|$ is the query complexity, and $\stmt, \wit$ denote the public statement and witness for the inner SNARK.}
\label{fig:protocol}
\end{figure*}

\paragraph{Remark on the SNARK relation.}
The SNARK proof $\pi$ attests to the validity of the relation $\mathcal{R}_{\mathsf{SNARK}}$ defined in Section~\ref{sec:relation-chain}.
The SNARK circuit verifies: (1)~constraint~\eqref{eq:codecom-check}, ensuring $\cm_x,\cm_y,\cm_z$ bind, via the selected code-commitment backend, to the full encodings of the same $\vect{x},\vect{y},\vect{z}$ used in the fold checks; (2)~constraints~\eqref{eq:fold-A}--\eqref{eq:eq-check}, the four fold/equality checks at each $j \in I$; and (3)~constraint~\eqref{eq:path-check}, consistency of all Merkle opening proofs with the committed roots.
The public statement $\stmt$ consists of the commitment roots $\{\cm_\chi\}_{\chi \in \mathcal{S}}$, the challenge $r$, and the query tuple $I$.
The witness $\wit$ consists of the intermediate vectors, the backend proofs/witnesses for their codeword commitments, and the opened codeword columns with their authentication paths.

\paragraph{Non-interactive variant.}
The protocol as described is a three-round public-coin interactive protocol.
By applying the Fiat--Shamir heuristic~\cite{fiat1987prove} in the random oracle model, it can be made non-interactive using domain-separated transcript hashes:
\[
  r \gets \Hash(\textsf{FMM-r},\para,\cm_A,\cm_B,\cm_C),
\]
and
\[
  I \gets \textsf{Sample}_{t}\!\left(
    \Hash(\textsf{FMM-I},\para,\cm_A,\cm_B,\cm_C,r,\cm_x,\cm_y,\cm_z)
  \right).
\]
The resulting non-interactive proof consists of the commitment roots, the
transcript-derived challenges, and the SNARK proof $\pi$; backend-specific
opening, linking, or backend code-commitment material is either carried as part of $\pi$ or
as auxiliary proof data verified by the corresponding backend.

\paragraph{Security of the non-interactive variant.}
A potential concern with the Fiat--Shamir transform for proximity-based protocols is a \emph{grinding attack}: a malicious prover could try many transcript prefixes until the derived challenge $r$ or query tuple $I$ is favorable.
The random-oracle theorem accounts for this by the standard
$\epsilon_{\mathsf{FS}}(q_H)$ loss in Theorem~\ref{thm:security}.
For intuition, consider grinding only the intermediate roots after $r$ is fixed.
Suppose the adversary has committed to $(\cm_A, \cm_B, \cm_C)$ (and hence to fixed encoded matrices) and, after receiving $r$, is searching for $(\cm_x, \cm_y, \cm_z)$ such that the transcript-derived query tuple $I$ avoids the disagreement set $D$ of some fold or equality check.
If any of the four checks fails globally, then $|D| \ge \delta n$ by the minimum-distance property.
For a single random query $j \in [n]$, the probability of $j \notin D$ is at most $1 - \delta$; for $t$ independent queries it is $(1-\delta)^t$.
Choosing $t$ as in Section~\ref{sec:security} makes the union-bounded proximity term $4(1-\delta)^t$ at most $2^{-\lambda}$.
Each trial of the grinding attack produces an independently random $I$ in the random oracle model because the query hash binds the full transcript, so the probability that $q$ grinding attempts yield a favorable $I$ is at most $q \cdot 4(1-\delta)^t$.
For any $\ppt$ adversary, $q$ is polynomially bounded in $\lambda$, giving a total success probability of $\mathsf{negl}(\lambda)$.
Hence the non-interactive protocol retains the soundness bound of
Theorem~\ref{thm:security} in the random oracle model, plus the stated
Fiat--Shamir loss.

\paragraph{Zero-knowledge extension.}
The protocol as stated is a \emph{verifiable computation} protocol and does not conceal the inputs $\mat{A}, \mat{B}, \mat{C}$ from the verifier.
Zero-knowledge can be achieved straightforwardly by two modifications:
(i) replace the Merkle tree commitments with a \emph{hiding} commitment scheme (e.g., salted hashes or Pedersen vector commitments~\cite{pedersen1991non}) so that the committed encoded matrices reveal no information about $\mat{A}, \mat{B}, \mat{C}$; and
(ii) employ a zero-knowledge SNARK backend for the inner proof $\pi$.
Masking the intermediate vectors $\vect{x}, \vect{y}, \vect{z}$ requires adding fresh random blinding vectors in Phase~3 and including the corresponding decommitment randomness in the SNARK witness.
We omit a full ZK construction and proof here; hiding the sampled openings,
the backend code-commitment material, and any linking material requires a separate
simulation argument.


\subsection{Complexity Analysis}
\label{sec:complexity}

We analyze the computational and communication costs of $\protocol$.

\paragraph{Prover computation.}
The prover's dominant costs are:
\begin{enumerate}[label=(\roman*)]
  \item \emph{Matrix multiplication}: Computing $\mat{C} = \mat{A}\mat{B}$ requires $\mathcal{O}(k^3)$ field operations (or $\mathcal{O}(k^\omega)$ with fast algorithms). This is performed natively outside the SNARK circuit and is amenable to GPU acceleration.
  \item \emph{Encoding}: $3k$ rows (one per matrix row, three matrices) each encoded into length-$n$ codewords costs $\mathcal{O}(kn)$ per matrix with linear-time codes.
  \item \emph{Vector computation}: Computing $\vect{x}, \vect{y}, \vect{z}$ costs $\mathcal{O}(k^2)$ in total.
  \item \emph{SNARK proving}: The circuit size is $\mathcal{O}(t \cdot k + t \cdot \log n)$ for the fold and Merkle checks, plus the backend codeword-validity term $E_{\mathsf{cc}}(k)$ (see below). The SNARK prover time scales with the circuit size.
\end{enumerate}

\paragraph{Verifier computation.}
The verifier performs SNARK verification together with any commitment/linking
checks required by the chosen backend.
For a pairing-based SNARK whose public input consists only of the roots,
challenge, and query indices, the core SNARK verification is $\mathcal{O}(1)$.
For transparent or linking-heavy backends, verification may scale with the
public input or linking proof size; the concrete behavior is reported in
Section~\ref{sec:implementation}.

\paragraph{Communication.}
The communication depends on the commitment layer.
With a fully succinct vector-commitment backend, the proof consists of a
constant number of commitment digests plus the SNARK proof, giving
$\mathcal{O}(\lambda)$ communication.
With the Merkle/CPLink implementation evaluated in this paper, the proof also
contains authentication, linking, and code-commitment backend material for the
queried encoded columns.
The sampled-column material scales as $\mathcal{O}(tk + t\log n)$ field/hash
elements for fixed security parameter; any additional backend material is
backend dependent and accounted for in the concrete proof sizes.
This trade-off explains the megabyte-scale proof sizes observed in
Section~\ref{sec:eval-verify}.

\paragraph{Circuit size comparison.}
Table~\ref{tab:cost-summary} summarizes the asymptotic costs.
Na\"ive arithmetization of $\mat{A}\mat{B} = \mat{C}$ requires $\mathcal{O}(k^3)$ constraints.
Freivalds' algorithm alone reduces this to $\mathcal{O}(k^2)$.
Our approach achieves $\mathcal{O}(tk + t\log n + E_{\mathsf{cc}}(k))$:
the first two terms cover the fold and Merkle checks, while
$E_{\mathsf{cc}}(k)$ denotes the cost of verifying the code-commitment backend
relation for the intermediate vector commitments.
For dense general codes, $E_{\mathsf{cc}}(k)=\mathcal{O}(k(n-k))$, which is
$\mathcal{O}(k^2)$ at constant rate.
For linear-time encodable codes of rate $\rho$ and distance $\delta$,
$E_{\mathsf{cc}}(k)=\mathcal{O}(k)$ when the selected code-commitment backend
has a linear-time verifier,
and setting
$t = \lceil(\lambda+2) / \log_2(1/(1-\delta))\rceil$ gives a total circuit size
of $\mathcal{O}(k)$ for fixed security parameter, with concrete sizes
discussed in Section~\ref{sec:implementation}.

\begin{table}[t]
\centering
\begin{tabular}{lcc}
\toprule
\textbf{Component} & \textbf{Prover} & \textbf{Verifier} \\
\midrule
Matrix multiplication & $\mathcal{O}(k^3)$ & --- \\
Encoding (row-wise) & $\mathcal{O}(kn)$ & --- \\
Merkle commitment & $\mathcal{O}(kn)$ & --- \\
SNARK circuit size & $\mathcal{O}(tk + t\log n + E_{\mathsf{cc}}(k))$ & --- \\
Core SNARK verification & --- & $\mathcal{O}(1)$ \\
\midrule
\textbf{Total} & $\mathcal{O}(k^3 + kn)$ & backend-dependent \\
\bottomrule
\end{tabular}
\caption{Asymptotic cost comparison. $k$: matrix dimension; $n$: codeword length; $t$: query count; $E_{\mathsf{cc}}(k)$: cost of the selected code-commitment backend relation.}
\label{tab:cost-summary}
\end{table}
